% !TEX root = ../Thesis.tex
% ============================================================
\chapter{Background} \label{chapter:background}
% ============================================================

\section{Knowledge Graphs and Graph Databases} \label{sec:kg-neo4j}

A knowledge graph represents information as a set of entities and the relations
between them, rather than as rows in a table or paragraphs in a document.
Concretely, it consists of nodes, each standing for one entity such as a person,
a drug, or a diagnosis, connected by directed, typed edges that state how two
entities relate, for instance that a particular drug treats a particular
condition. This structure matters for two reasons that recur throughout this
thesis. First, it is explicit: every fact is a node or an edge that can be looked
up, corrected, or removed on its own, rather than being buried inside a block of
free text. Second, it suits questions that require connecting several facts,
since answering such a question becomes a matter of following edges from one node
to the next rather than searching through unrelated documents.

Knowledge graphs are stored and queried through graph databases, and one of the
most widely used is Neo4j~\cite{neo4j2026}. Neo4j represents data as a property
graph, meaning that both nodes and edges can carry arbitrary key--value
properties in addition to their type, and it is queried with
Cypher~\cite{francis2018cypher}, a query language built for describing graph
patterns rather than table joins. A query such as
\texttt{MATCH (d:Drug)-[:TREATS]->(c:Condition \{name: "X"\}) RETURN d} reads
close to the pattern it looks for: find drug nodes connected by a \texttt{TREATS}
edge to a condition node named X. Neo4j is the database underlying both SnapQuery
and the retrieval method developed here, so the graph this thesis works with is,
at the storage level, a Neo4j property graph, even though the schema that defines
its structure is expressed in a different formalism
(Section~\ref{subsec:sphn-schema}).

\section{Retrieval-Augmented Generation} \label{sec:rag-background}

Large language models generate fluent text, but the knowledge stored in their
weights during training is fixed and general purpose, which leaves them prone to
hallucination: stating something confidently that is simply wrong.
Retrieval-Augmented Generation (RAG), introduced by Lewis et
al.~\cite{lewis2020rag}, addresses this by retrieving relevant evidence from an
external knowledge source at query time and conditioning the model's output on
that evidence, rather than relying solely on what the model already encodes
internally. In its original form, the knowledge source is a large collection of
text passages, and retrieval means picking out the passages whose text is closest
to the query.

This works well for many tasks, but it has one limitation that matters here: a
flat collection of passages carries no relational structure. Two passages
connected through some intermediate concept, say two drugs that each interact
with the same enzyme, are treated as entirely independent by the retriever, and
the retrieval step cannot follow a chain of relations from one to the other. That
becomes an obstacle whenever an answer depends on connecting several facts, which
is exactly what knowledge graph question answering requires. A question such as
``which drugs treat a disease that shares a symptom with condition X?'' cannot be
answered from isolated passages, because the reasoning path runs through the
graph rather than sitting inside any single one of them.

\section{From Text Retrieval to Graph Retrieval} \label{sec:graph-retrieval}

Several lines of work adapt retrieval to graph-structured knowledge, each
addressing that limitation differently. KAPING~\cite{baek2023kaping} retrieves
the top-$k$ triples ranked most similar to the query by embedding similarity, a
measure of how close two pieces of text are in meaning computed as the distance
between their vector representations, and prepends them as flat text to the
prompt. This is fast and simple, but it inherits the connectivity blindness of
text-based RAG, since triples are still selected independently with no guarantee
that they form a coherent chain. Think-on-Graph~\cite{sun2024thinkongraph} takes
the opposite approach: it performs an iterative search over the graph, using the
LLM itself to decide which relation to follow at each step. The paths it produces
are more coherent, but the method needs several LLM calls per query, which becomes
expensive at scale. StructGPT~\cite{jiang2023structgpt} and
MindMap~\cite{wen2023mindmap} construct explicit reasoning chains over a
knowledge graph before generation, with MindMap applied specifically to
biomedical problems.

A related idea, Microsoft's GraphRAG~\cite{edge2024graphrag}, builds
community-level summaries over an LLM-constructed graph, which suits broad,
corpus-level questions better than the entity-specific questions studied here.
QA-GNN~\cite{yasunaga2021qagnn} pairs a language model directly with a graph
neural network; Section~\ref{sec:gnn-llm} returns to it once graph neural
networks have been introduced. What these methods share is the insight that
incorporating graph structure into retrieval improves on flat text retrieval, and
what they also share is that they were designed for small or synthetic graphs
rather than for the large, dense graphs typical of hospital data.

\section{Subgraph Selection Approaches} \label{sec:subgraph-selection-approaches}

The methods above retrieve individual pieces of a graph, but say little about how
those pieces should be combined into the connected region eventually handed to
the LLM. What separates the available families in practice is how each one
handles connectivity, and Table~\ref{tab:selection-families} summarises where
they stand.

The simplest strategy, top-$k$ selection, scores every node and edge by
similarity to the query and keeps the $k$ highest-scoring ones, with $k$ fixed in
advance. Because each item is scored independently this is fast, and the size of
the result is known exactly, but there is no guarantee that the returned facts
relate to one another, which hurts both answer quality and interpretability.
Neighbourhood expansion methods add connectivity by collecting the $h$-hop
neighbourhood of the selected nodes, meaning every node reachable within $h$
steps. Connectivity is then guaranteed by construction, but size control is lost:
in a dense graph the neighbourhood grows by orders of magnitude per hop, and $h$
is an integer, so there is no setting available between too small and far too
large. Path-based methods instead connect the selected entities through the
graph, which also guarantees a connected result, but they can pull in long paths
through unrelated intermediate nodes and offer no direct handle on subgraph size.
LLM-guided traversal, as in Think-on-Graph, decides each step with a model call
and so pays a cost per question that the others do not.

PCST, described in the next section, is the only family in the table that offers
both a connectivity guarantee and a continuous size parameter, because it treats
the two as terms of a single objective rather than as separate stages.

\begin{table}[t]
\centering
\caption{Families of subgraph selection strategies. ``Size control'' asks whether
the strategy exposes a parameter that predictably bounds the size of the returned
subgraph on a \emph{dense} graph; this is where neighbourhood expansion and
path-based methods differ from the other three.}
\label{tab:selection-families}
\begin{tabular}{llll}
\toprule
Family & Connected result & Size control & Cost per question \\
\midrule
Top-$k$ similarity      & no  & exact, via $k$            & one encoding pass \\
Neighbourhood expansion & yes & coarse, via hops $h$      & one traversal \\
Path-based              & yes & none direct               & one or more traversals \\
LLM-guided traversal    & yes & none direct               & several LLM calls \\
PCST                    & yes & continuous, via edge cost & one optimisation \\
\bottomrule
\end{tabular}
\end{table}

\section{G-Retriever and Prize-Collecting Steiner Tree Selection} \label{sec:gretriever-pcst}

G-Retriever~\cite{he2024gretriever} is the system this thesis builds on. This
section describes what it does and why the formulation is attractive; Chapter~3
gives the mathematical treatment, including the exact prize and cost definitions,
the virtual-node transformation that lets edges carry prizes, and the
approximation guarantee of the solver used. The division is deliberate: the
argument of this chapter does not depend on the algebra, and the algebra is
easier to read once the argument is in place.

G-Retriever runs in four stages. First, every node and edge in the graph has its
text description encoded into a vector with a sentence
encoder~\cite{reimers2019sentencebert}, and the query is encoded the same way;
the nodes and edges whose vectors sit closest to the query's become the initial
candidates, exactly as in the similarity-based methods above. Second, rather than
returning those candidates directly, G-Retriever treats subgraph construction as
an optimisation problem: find a connected subgraph that keeps as much of the
total relevance, or ``prize'', of the candidate nodes and edges as possible while
paying a cost for every edge needed to connect them. This is the
Prize-Collecting Steiner Tree (PCST) problem~\cite{bienstock1993pcst}, a
classical combinatorial optimisation problem originally studied in network
design. Framing selection this way has a useful consequence: a node of only
moderate relevance can still be worth including if it lies on a short path
between two highly relevant nodes, because the prize gained by connecting them
can outweigh the cost of the extra edge. Such a node is invisible to any method
that scores nodes independently.

Third, once the connected subgraph has been found, its structure is encoded by a
graph neural network into a short vector summary, and fourth, that summary is
prepended to the LLM's input as an additional, learned token, often called a soft
prompt because it behaves like part of the prompt without being ordinary text. On
the general-purpose benchmarks G-Retriever was evaluated on, the pipeline cut the
number of tokens needed by up to 99\% compared with serialising the whole graph
as text, while also reducing hallucination~\cite{he2024gretriever}. Whether those
benefits carry over to a dense, specialised clinical graph is the central
empirical question of this thesis.

\section{Graph Neural Networks and LLMs on Graphs} \label{sec:gnn-llm}

The soft-prompt construction just described is one of two broad strategies for
combining a graph neural network (GNN), a network that computes a representation
for each node by repeatedly combining it with its neighbours'
representations~\cite{velickovic2018gat}, with a language model. It compresses
the graph into a fixed prefix before generation begins. The other strategy,
illustrated by QA-GNN~\cite{yasunaga2021qagnn}, lets the GNN and the language
model exchange information throughout encoding, so the two representations inform
each other as they are built rather than one being frozen before the other
starts. Applied to medical question answering, QA-GNN first extracts a small
question-specific subgraph and then runs a GNN over it jointly with the language
model.

QA-GNN matters here for two reasons beyond the architectural contrast. It is an
early demonstration that a question-specific subgraph combined with a graph
encoder improves question answering over text-only retrieval, and its two-stage
shape, isolate a subgraph first and encode it second, is the shape studied in this
thesis. Where this thesis differs is in where it puts its attention: QA-GNN and
G-Retriever both treat subgraph extraction as a preliminary and the encoder as
the contribution, whereas here the extraction step \emph{is} the object of study,
and the generation stage is held fixed so that differences in the final answer
can be attributed to what was retrieved. For a broader overview of how GNNs and
LLMs have been combined on graph tasks, see the survey by Jin et
al.~\cite{jin2023llmgraphsurvey}.

\section{The Clinical Setting: STCS, SPHN, and SnapQuery} \label{sec:clinical-kg}

\subsection{Curated Ontologies and Mined Knowledge Graphs} \label{subsec:curated-ontologies}

Clinical and biomedical knowledge graphs are built from two main kinds of source.
Curated ontologies encode expert-validated medical knowledge as structured
hierarchies and relational networks, with well-typed nodes and edges that can
generally be trusted for clinical applications. Two are used widely enough to
name here: SNOMED~CT (Systematized Nomenclature of Medicine Clinical
Terms)~\cite{snomedct}, a comprehensive clinical vocabulary whose concepts are
arranged in explicit is-a relationships rather than a flat list, and the UMLS
(Unified Medical Language System)~\cite{bodenreider2004umls}, which maps concepts
from many such vocabularies onto one another. Automatically mined
knowledge graphs, built from electronic health records~\cite{rotmensch2017learning},
cover emerging or institution-specific knowledge but may contain noise, since
they are extracted rather than curated.

The graph studied here is of a third kind, which is worth naming because it
changes what retrieval has to cope with. It is neither a curated ontology of
general medical knowledge nor a mined approximation of one, but a
schema-constrained record of \emph{what happened to particular patients}. Its
nodes are events and measurements rather than concepts, and general medical
knowledge enters it only indirectly, through the terminology codes those events
reference (Section~\ref{subsec:terminologies}). A retriever therefore cannot rely
on the semantic richness that makes ontology nodes easy to match against a
question, because most nodes in this graph are not concepts at all.

\subsection{The Swiss Transplant Cohort Study} \label{subsec:stcs-background}

The Swiss Transplant Cohort Study (STCS) is a nationwide, prospective cohort that
has enrolled essentially all solid organ transplant recipients in Switzerland
since 2008, across the country's transplant centers, including University
Hospital Basel~\cite{koller2013stcs}. It was established because organ
procurement in Switzerland was already well organised at the national level while
transplant outcomes were not being tracked systematically across centers. STCS
collects detailed, longitudinal clinical data on each patient, covering the
transplantation itself, subsequent complications such as infections and rejection
episodes, and long-term follow-up, and has since supported several hundred
research projects in transplantation medicine. The data this thesis works with is
the STCS data held at University Hospital Basel, represented as a knowledge graph
rather than as the flat tables the cohort was originally collected in.

That change of representation is not cosmetic. A cohort table has one row per
patient per visit, and its columns are the questions the cohort was designed to
answer. A graph has no privileged shape, so a question nobody anticipated can
still be expressed. That is the reason for building the graph, and it is also the
reason retrieval over it is hard: nothing in the structure marks which nodes
matter for a given question.

\subsection{The Swiss Personalized Health Network (SPHN) Schema} \label{subsec:sphn-schema}

Turning STCS data into a usable knowledge graph requires a schema that fixes what
kinds of entities and relations may appear in it, and in Switzerland that schema
is defined nationally rather than by each hospital independently. The Swiss
Personalized Health Network (SPHN) schema~\cite{sphnschema} is written in the
Resource Description Framework (RDF), a W3C standard that represents data as
subject--predicate--object triples and lets a schema fix the permitted concept
classes and the properties linking them~\cite{rdfschema}. Aligning to a national
standard rather than to a one-off local schema keeps a pipeline built on it
portable across hospital sites, which matters because the STCS cohort spans every
transplant center in the country.

Two properties of the schema shape everything that follows, and both were
confirmed against the live graph during this project.

\paragraph{Reification.} SPHN models values, units, codes and body sites as
\emph{their own concept classes} rather than as attributes of the clinical event
that uses them. In the property-graph rendering, each RDF class becomes a node
label, each object property becomes a typed relationship, and only datatype
properties become node properties. A single laboratory result is therefore not
one node but a small subgraph: an observation node linked to a quantity node,
which is linked to a unit node, while the observation is also linked to a test
type, which is linked to a code, which carries the human-readable description.
One clinical fact spans on the order of five to seven nodes, which is the direct
explanation for the roughly 18{,}000 nodes per patient reported in
Section~\ref{sec:problem-statement}.

\paragraph{Text lives on shared nodes, not on event nodes.} Because descriptions
sit on code nodes, and code nodes are shared by every event that references them,
the text available to an embedding-based retriever is attached to the wrong
objects. The event a question is actually about carries a timestamp and an
identifier; the words a question can match live one or two hops away, on a node
shared with thousands of other events. This is what produces the
sparse-and-shared-text property reported in
Section~\ref{sec:problem-statement}, and it is why Chapter~4 introduces a graph
preparation step that folds shared terminology text into the events that
reference it before any embedding is computed.

% TODO (confirm with the STCS Data Center): the label carrying diagnoses in this
% instance is BilledDiagnosis. Confirm that billed diagnoses are the right
% clinical object for the questions users actually ask, rather than a separate
% clinically-recorded diagnosis label. If not, the diagnosis queries change but
% nothing else in the pipeline does.
The concept classes relevant to this thesis are the subject identifier, the
administrative case that groups one hospital stay, laboratory observations with
their test types, diagnoses, and drug administration events with the drugs they
administer. Alongside these sit provenance classes recording which source system
and data release each fact came from. Those provenance nodes are the ones whose
degree was measured in Section~\ref{sec:problem-statement}: they are correct,
required by the schema, and ruinous for connectivity-based retrieval, because they
place almost every pair of nodes within a few hops of one another.

\subsection{Clinical Terminologies and the Multilingual Problem} \label{subsec:terminologies}

SPHN does not invent a vocabulary of clinical concepts but reuses established
\emph{terminologies}: published catalogues that assign a stable code to each
clinical concept, so that the same test or diagnosis is written the same way in
every hospital that adopts them. A node in the graph carries the code; the words
belonging to that code live in the catalogue. Five such catalogues appear
together in the graph studied here.

Laboratory tests are coded with LOINC (Logical Observation Identifiers Names and
Codes)~\cite{loinc}, the international standard for identifying laboratory
measurements, whose English descriptions are present for every one of the 544
distinct test codes in the cohort. Diagnoses are coded with ICD-10-GM, the German
national modification of the World Health Organization's International
Classification of Diseases, tenth revision~\cite{icd10}. It is versioned by year,
and 2{,}921 distinct codes appear across thirteen annual versions. Drugs are
coded with ATC (Anatomical Therapeutic Chemical)~\cite{atc}, a World Health
Organization classification that groups a substance by the organ it acts on and
its therapeutic properties. Units of measurement use UCUM (Unified Code for Units
of Measure), a machine-readable notation in which, for instance, millimoles per
litre is written \texttt{mmol/L}. A small number of concepts use
SNOMED~CT, introduced in Section~\ref{subsec:curated-ontologies}.

This mixture creates a problem that the general-domain literature does not have
to confront, and it is easy to miss because it announces itself in no error
message. The descriptions carried by the graph are in different languages
depending on which terminology a node belongs to: LOINC descriptions are English,
ICD-10-GM descriptions are German. G-Retriever's sentence encoder,
\texttt{all-roberta-large-v1}, is English-only. Embedding the labels the graph
happens to carry therefore under-ranks the entire diagnosis portion of the graph
relative to the laboratory portion, for a reason that has nothing to do with
clinical relevance. Because the distortion correlates with entity type rather
than with question difficulty, it does not appear as noise in an aggregate recall
figure; it appears as a consistent and plausible-looking preference for
laboratory evidence. A second, compounding gap is that 1{,}949 of the 2{,}921
diagnosis codes carry no description in the graph at all, leaving nodes that
similarity search cannot reach under any wording of any question.

Both problems have the same solution, and it is not translation. A terminology
code is language-independent by construction: \texttt{N18.5} denotes the same
concept whichever language its label is written in. Resolving codes against their
published catalogues therefore replaces a translation problem with a lookup, and
a lookup is deterministic, auditable and reproducible in a way that machine
translation of clinical shorthand is not. It also fills the missing descriptions,
because a catalogue entry exists whether or not the local graph stored a label.
Chapter~4 specifies the procedure, including the fallback tiers used when no
catalogue entry is found and the reporting of how many nodes each tier accounts
for, and Chapter~6 reports what difference it makes. The point to carry forward
is that terminology resolution is not preprocessing hygiene in this setting but a
precondition for the retrieval signal to mean anything, and that it is also the
mechanism by which results obtained on a German-coded cohort stay meaningful for
a cohort coded in another language.

\subsection{SnapQuery: The System Already in Use at STCS} \label{subsec:snapquery-background}

SnapQuery is the assistant that clinicians and researchers at the STCS Data
Center, University Hospital Basel, use today to ask questions about hospital data
in plain language~\cite{snapquery2026}. It sits as an access layer between its
users and the hospital's data backends: given a question, it has a language model
write a single read-only database query, executes that query, and returns a
structured answer bundling a natural-language response, the generated query
itself, and the matching rows. Its purpose is to remove the data manager as a
mandatory intermediary for routine access, not to replace the data infrastructure
it queries. SnapQuery supports two independent paths, a graph path against three
hospital Neo4j knowledge graphs (University Hospital Basel, Inselspital Bern and
CHUV Lausanne) and a tabular path against a roughly 230-table PostgreSQL clinical
warehouse at Basel. Only the graph path is relevant here, since the STCS cohort
studied in this thesis is one of the graph-backed centers.

The graph path is a LangGraph agent: an assistant component performs the LLM
reasoning and decides which tool to call, a tool component executes that action,
and the two alternate until an answer is ready. Its system prompt embeds the
session center's live Neo4j schema together with a slot-filling policy that asks
at most two clarifying questions per turn to establish details such as the time
window, cohort or output shape, and it generates and executes a query only once
the user has explicitly confirmed what is about to run. Schema lookup and Cypher
generation are delegated to a self-hosted planner model, Qwen2.5-14B-Instruct,
which reasons over the schema returned by a dedicated lookup tool and produces
Cypher through a retrieval-augmented chain that validates the query before it
runs, with preview rows capped between one and fifty. One design difference is
worth keeping in mind: unlike the tabular path, the graph agent uses no vector
store for schema retrieval, judging relevance by having the planner model read
the full live schema, whereas the method developed in this thesis is
embedding-based throughout. The comparison in Chapter~6 is therefore between a
schema-reading query generator and a text-similarity retriever, not between two
variants of one mechanism.

What matters most for this thesis is what SnapQuery ultimately returns: one
generated Cypher query and the rows it produces, with no explicit subgraph of
intermediate nodes or edges. Chapter~4 treats the entities touched by that result
set as SnapQuery's induced subgraph, which is what makes it possible to score the
two systems on a single metric. One consequence of that choice, that SnapQuery
yields a single point rather than a curve because it exposes no size parameter to
sweep, is stated there as a limitation of the comparison rather than hidden
inside it.

\subsection{LLM Prompting over Clinical Knowledge Graphs} \label{subsec:llm-clinical-prompting}

Recent work has begun combining clinical knowledge graphs with LLM prompting for
decision support. Soman et al.~\cite{soman2024biomedical} use a biomedical
knowledge graph to construct optimised prompts for LLMs, and
MindMap~\cite{wen2023mindmap} builds explicit reasoning maps over a medical
knowledge graph before generation. Both draw on curated biomedical knowledge
rather than on patient-level hospital records, which is exactly the distinction
drawn in Section~\ref{subsec:curated-ontologies}. To the best of our knowledge,
no existing work applies connectivity-constrained, PCST-style subgraph selection
to a patient-level clinical knowledge graph, evaluates such selection against a
deployed clinical query system, or examines the effect of intent and schema
signals within such a retrieval framework.

\section{Summary: What the Surveyed Methods Assume} \label{sec:background-summary}

The methods reviewed in this chapter share three assumptions, and the graph
measured in Section~\ref{sec:problem-statement} violates all three. Naming them
here converts the survey above into three specific, testable predictions, which
Chapter~3 formalises and Chapter~6 tests.

The first assumption is that node text is discriminative: that similarity to a
question can tell one node from another. Under reification and shared terminology
text (Sections~\ref{subsec:sphn-schema} and~\ref{subsec:terminologies}) thousands
of nodes carry identical text, so the ranking within a tied group is decided by
whatever order the implementation happens to iterate in rather than by relevance.
The prediction is that similarity can identify which \emph{kind} of fact a
question concerns but not \emph{which instance}, and that instance selection must
therefore come from structure or from time.

The second assumption is that a relation type is informative about relevance.
G-Retriever assigns prizes to edges by ranking relation text against the
question, dividing a fixed budget among the edges tied at each rank. That is
reasonable when a graph has many relation types each used a few times, as in the
open-domain benchmarks. It fails when a graph has few relation types each used
millions of times, because the per-edge share of the budget becomes negligible
and, through the cost-adjustment step that keeps costs below the largest prize,
drags the effective edge cost towards zero with it. The prediction is that the
size dial stops working and the method returns everything it is given.

The third assumption is that connectivity is a scarce resource worth paying for.
That holds in a sparse graph of clinical concepts. It does not hold once
provenance nodes of very high degree are present, because then almost any two
nodes are already close together and a connectivity requirement costs nothing to
satisfy. The prediction is that connectivity becomes meaningful only after such
nodes are excluded, and that the effect of excluding them shows up in what the
retrieved subgraph contains rather than in how large it is.

None of the three is a flaw in G-Retriever, which was designed for and validated
on a different kind of graph. They are the price of the transfer, and the reason
this thesis measures it rather than assuming it. Chapter~4 turns each prediction
into a methodological choice that can be enabled and disabled independently, so
that the contribution of each can be reported on its own.
